\section{TSK mactime}

Einfache Timeline aus Sleuth Kit Body File:

\begin{lstlisting}[language=bash]
# Body File erzeugen (alle Dateien mit Timestamps)
fls -r -m "/" -o OFFSET image.E01 > body.txt

# Timeline generieren
mactime -b body.txt > timeline.csv
\end{lstlisting}

\begin{description}
    \item[MAC-Timestamps] Modified, Accessed, Changed
    \item[Body File Format] Enthält Inode, Permissions, UID, GID, Größe und alle Timestamps
\end{description}

\section{Plaso / log2timeline}

Mächtigere Super-Timeline — analysiert viele Artefakttypen über reine Datei-Timestamps hinaus:

\begin{lstlisting}[language=bash]
# Timeline aus gemounteter Partition erstellen
log2timeline.py /tmp/timeline.plaso /mnt

# Nur bestimmte Parser verwenden (schneller)
log2timeline.py --parsers "chrome_history,firefox_history" \
    /tmp/browser.plaso /mnt

# Als CSV exportieren
psort.py -o l2tcsv -w /tmp/timeline.csv /tmp/timeline.plaso
\end{lstlisting}

Verfügbare Parser für Web-History:
\begin{itemize}
    \item \texttt{chrome\_history}, \texttt{firefox\_history}, \texttt{msedge\_history}
    \item \texttt{ie\_webcache}, \texttt{safari\_history}, \texttt{opera\_history}
\end{itemize}

\section{Browser-History direkt per SQLite}

Alternative: Datenbanken direkt extrahieren und abfragen.

\subsection{Chrome}

\begin{lstlisting}[language=bash]
fls -pr /dev/loopXpN | grep -i 'History$' | grep -i chrome
icat /dev/loopXpN <INODE> > /tmp/chrome_history.db
sqlite3 /tmp/chrome_history.db \
  "SELECT url, title, visit_count, last_visit_time
   FROM urls ORDER BY last_visit_time DESC LIMIT 50;"
\end{lstlisting}

\subsection{Firefox}

\begin{lstlisting}[language=bash]
fls -pr /dev/loopXpN | grep -i 'places.sqlite'
icat /dev/loopXpN <INODE> > /tmp/places.sqlite
sqlite3 /tmp/places.sqlite \
  "SELECT url, title, visit_count, last_visit_date
   FROM moz_places ORDER BY last_visit_date DESC LIMIT 50;"
\end{lstlisting}

\subsection{Edge (ESE-Datenbank)}

\begin{lstlisting}[language=bash]
fls -pr /dev/loopXpN | grep -i 'WebCacheV01'
icat /dev/loopXpN <INODE> > /tmp/WebCacheV01.dat
esedbexport -t /tmp/edge_cache /tmp/WebCacheV01.dat
\end{lstlisting}
